\begin{trapbox}

	\begin{description}[style=nextline, leftmargin=2.8em, labelsep=0.5em, itemsep=0.35em]
		\item[\textbf{\textcolor{CTrap}{易错点一（充要混淆）}}]
			认为 $f'(x_0)=0$ 则 $x_0$ 一定是极值点.
		\item[\textbf{\textcolor{CTrap}{正确理解（必要不充分）}}]
			$f'(x_0)=0$ 只是必要条件，不能直接判定极值，需验证左右导数符号或二阶导数.
		\item[\textbf{\textcolor{CTrap}{错因分析（概念混淆）}}]
			混淆必要条件和充分条件，忽略了反例 $f(x)=x^3$.
	\end{description}

	\medskip
	\begin{description}[style=nextline, leftmargin=2.8em, labelsep=0.5em, itemsep=0.35em]
		\item[\textbf{\textcolor{CTrap}{易错点二（验证缺失）}}]
			利用极值点必要条件求参数后，直接认为参数满足条件，不验证该点是否为极值点.
		\item[\textbf{\textcolor{CTrap}{正确理解（参数检验）}}]
			求出参数后必须代回验证，确认该点为极值点（否则可能为 $x^3$ 型驻点）.
		\item[\textbf{\textcolor{CTrap}{错因分析（步骤缺失）}}]
			解方程仅得到驻点条件，但驻点未必是极值；忽略充分性验证.
	\end{description}

	\medskip
	\begin{description}[style=nextline, leftmargin=2.8em, labelsep=0.5em, itemsep=0.35em]
		\item[\textbf{\textcolor{CTrap}{易错点三（前提忽略）}}]
			认为所有极值点处导数都必须为零，包括不可导点.
		\item[\textbf{\textcolor{CTrap}{正确理解（可导前提）}}]
			该结论仅对可导函数成立；在不可导点也可能出现极值，此时导数不存在.
		\item[\textbf{\textcolor{CTrap}{错因分析（条件遗忘）}}]
			忽略“可导”这一条件，将结论推广到所有函数.
	\end{description}

\end{trapbox}
